\section{Theoretical Mixing Time Bounds}
\label{sec:bounds}

\subsection{Spectral Gap Framework}

For a reversible Markov chain with transition matrix $P$ and stationary
distribution $\pi$, the spectral gap is:
\[
  \sgap = 1 - \lambda_2,
\]
where $\lambda_2$ is the second-largest eigenvalue of $P$.
The total variation mixing time satisfies:
\[
  \tmix(\epsilon) \leq \frac{\log(|\mathcal{P}| / \epsilon)}{\sgap}.
\]

To bound $\tmix$, we need a lower bound on $\sgap$.

\subsection{Conductance and the Cheeger Inequality}

The conductance of a reversible Markov chain is:
\[
  \Phi = \min_{S \subset \mathcal{P},\, 0 < \pi(S) \leq 1/2}
  \frac{\sum_{x \in S, y \notin S} \pi(x) P(x,y)}{\pi(S)}.
\]
The Cheeger inequality~\citep{sinclair1989} gives:
\[
  \frac{\Phi^2}{2} \leq \sgap \leq 2\Phi.
\]

A lower bound on $\Phi$ gives a lower bound on $\sgap$.

\subsection{The Planar Separator Bottleneck}

The \recom{} chain on plan space inherits its bottleneck from the
bottleneck in the underlying census-tract graph $G$.
Specifically, for the chain to move between two plan classes separated
by a bottleneck in $G$, all tracts in the bottleneck must change district
assignment.
The Lipton-Tarjan planar separator theorem~\citep{lipton1979} guarantees
that every planar graph on $n$ vertices has a separator of size $O(\sqrt{n})$.

\begin{theorem}[Mixing Time Upper Bound]
\label{thm:mixing}
For the lazy reversible \recom{} chain on a planar census-tract graph
$G$ with $n$ vertices and $k$ districts, the spectral gap satisfies:
\[
  \sgap = \Omega\!\left(\frac{1}{n^2}\right),
\]
and consequently:
\[
  \tmix(0.01) = O\!\left(n^3 \log n\right).
\]
\end{theorem}

\begin{proof}[Proof sketch]
The canonical path argument~\citep{diaconis1996} assigns to each pair
$(\pi, \pi')$ a path in $\mathcal{P}$ connecting them.
The bottleneck of the path argument is bounded by the number of paths
through each plan, which is bounded by the product of the separator size
and the number of plan pairs that must route through the separator.

For a planar graph with separator $S$ of size $O(\sqrt{n})$:
the number of plans that can be ``on either side'' of the separator
is at most $|\mathcal{P}| / |\mathcal{P}^*|$ where $\mathcal{P}^*$ is
the set of plans consistent with the separator assignment.
Standard bottleneck arguments give $\Phi = \Omega(1/n)$, and hence
$\sgap = \Omega(\Phi^2) = \Omega(1/n^2)$.

The mixing time bound follows from:
$\tmix = O(\log(|\mathcal{P}|) / \sgap) = O(n \log n \cdot n^2) = O(n^3 \log n)$.
\end{proof}

\begin{remark}
The tighter $O(n^2 \log n)$ bound would require a direct path congestion
argument showing that the congestion of the canonical paths is $O(n/\log n)$
rather than $O(n)$.  This requires path congestion analysis beyond the scope
of this paper; see \citet[Theorem~13.1]{levin2009markov} for the relevant
machinery.  Reducing the exponent from 3 to 2 does not affect the empirical
conclusions: both bounds are many orders of magnitude above observed mixing
times (see Table~\ref{tab:mixing-bounds}), confirming that theory-practice
gap analysis must proceed empirically.
\end{remark}

\subsection{Numerical Implications}

For six study states:

\begin{table}[h]
\centering
\caption{Theoretical mixing time upper bounds $O(n^3 \log n)$ and
         empirical $\rhat < 1.1$ thresholds. The gap between theory and
         practice spans 3-4 orders of magnitude.}
\label{tab:mixing-bounds}
\begin{tabular}{lcccc}
\toprule
State & $n$ (tracts) & Theoretical $\tmix$ & Empirical $\hat{t}_{\rm mix}$ & Ratio \\
\midrule
NC & 2{,}195 & $\sim$10$^8$ & 2{,}000  & $\sim$50{,}000$\times$ \\
WI & 1{,}406 & $\sim$4$\times$10$^7$ & 1{,}500  & $\sim$27{,}000$\times$ \\
GA & 1{,}969 & $\sim$8$\times$10$^7$ & 2{,}500  & $\sim$32{,}000$\times$ \\
PA & 3{,}218 & $\sim$3$\times$10$^8$ & 3{,}000  & $\sim$100{,}000$\times$ \\
TX & 5{,}265 & $\sim$7$\times$10$^8$ & 10{,}000 & $\sim$70{,}000$\times$ \\
CA & 8{,}057 & $\sim$2$\times$10$^9$ & 20{,}000 & $\sim$100{,}000$\times$ \\
\bottomrule
\end{tabular}
\end{table}

The theoretical bounds are worst-case; the empirical mixing times reflect
concentration near compact plans.
The ratio is consistent across states (50{,}000--100{,}000$\times$),
suggesting a systematic concentration effect rather than coincidence.

\subsection{Interpretation of the Gap}

The stationary distribution is approximately uniform over $\mathcal{P}$.
However, the high-compactness region of $\mathcal{P}$---the plans that
METIS would generate---is a small fraction of the total state space but
captures most of the stationary measure (in a compactness-weighted chain)
or is rapidly visited (in a uniform chain starting from a compact plan).

For a chain starting from a compact plan, the effective mixing time
within the high-compactness region is much smaller than the worst-case
mixing time over all of $\mathcal{P}$.
This explains the $10^4$--$10^5\times$ gap between theory and practice.
